\PassOptionsToPackage{table}{xcolor}
\PassOptionsToPackage{dvipsnames}{xcolor}
% \documentclass[pdf, aspectratio=169, xcolor=dvipsnames]{beamer}
\documentclass[pdf, aspectratio=169]{beamer}\usepackage[]{graphicx}\usepackage[]{color}
% maxwidth is the original width if it is less than linewidth
% otherwise use linewidth (to make sure the graphics do not exceed the margin)
\makeatletter
\def\maxwidth{ %
  \ifdim\Gin@nat@width>\linewidth
    \linewidth
  \else
    \Gin@nat@width
  \fi
}
\makeatother

\definecolor{fgcolor}{rgb}{0.345, 0.345, 0.345}
\newcommand{\hlnum}[1]{\textcolor[rgb]{0.686,0.059,0.569}{#1}}%
\newcommand{\hlstr}[1]{\textcolor[rgb]{0.192,0.494,0.8}{#1}}%
\newcommand{\hlcom}[1]{\textcolor[rgb]{0.678,0.584,0.686}{\textit{#1}}}%
\newcommand{\hlopt}[1]{\textcolor[rgb]{0,0,0}{#1}}%
\newcommand{\hlstd}[1]{\textcolor[rgb]{0.345,0.345,0.345}{#1}}%
\newcommand{\hlkwa}[1]{\textcolor[rgb]{0.161,0.373,0.58}{\textbf{#1}}}%
\newcommand{\hlkwb}[1]{\textcolor[rgb]{0.69,0.353,0.396}{#1}}%
\newcommand{\hlkwc}[1]{\textcolor[rgb]{0.333,0.667,0.333}{#1}}%
\newcommand{\hlkwd}[1]{\textcolor[rgb]{0.737,0.353,0.396}{\textbf{#1}}}%
\let\hlipl\hlkwb

\usepackage{framed}
\makeatletter
\newenvironment{kframe}{%
 \def\at@end@of@kframe{}%
 \ifinner\ifhmode%
  \def\at@end@of@kframe{\end{minipage}}%
  \begin{minipage}{\columnwidth}%
 \fi\fi%
 \def\FrameCommand##1{\hskip\@totalleftmargin \hskip-\fboxsep
 \colorbox{shadecolor}{##1}\hskip-\fboxsep
     % There is no \\@totalrightmargin, so:
     \hskip-\linewidth \hskip-\@totalleftmargin \hskip\columnwidth}%
 \MakeFramed {\advance\hsize-\width
   \@totalleftmargin\z@ \linewidth\hsize
   \@setminipage}}%
 {\par\unskip\endMakeFramed%
 \at@end@of@kframe}
\makeatother

\definecolor{shadecolor}{rgb}{.97, .97, .97}
\definecolor{messagecolor}{rgb}{0, 0, 0}
\definecolor{warningcolor}{rgb}{1, 0, 1}
\definecolor{errorcolor}{rgb}{1, 0, 0}
\newenvironment{knitrout}{}{} % an empty environment to be redefined in TeX

\usepackage{alltt}
% \documentclass[pdf, aspectratio=1610]{beamer}
% \setbeamersize{text margin left=2cm, text margin right=2cm}
\setbeamersize{text margin left=1.5cm, text margin right=1.5cm}
% \documentclass[pdf]{beamer}
% \usefonttheme[onlymath]{serif} % For serif font in math mode

\usepackage[italian]{babel} %La data compare in italiano e la divisione in sillabe diventa quella italiana
\usepackage[utf8]{inputenc} %Permette di scrivere le lettere accentate
\usepackage[T1]{fontenc}
\usepackage{hyperref} %Permette di inserire i collegamenti ipertestuali
\usepackage{url}			        %Permette di inserire i link urle
\usepackage{amsmath, amssymb}
\usepackage{multirow} % Per usare il multirow nei tabular
% \usepackage{enumitem} % For itemize options
\usepackage{fontawesome} % For emoji and custom symbols
\usepackage{stackengine} % For overlapping characters with \bottominset

% For itemize without spacing
\usepackage{tabularx}
\newcommand{\myitem}[1]{%
  \begin{tabularx}{\linewidth}{@{}l@{}>{\raggedright}X@{}}
  \usebeamercolor[fg]{structure}$\bullet$~~ & #1
  \end{tabularx}%
}   
\newenvironment{myitemize}{\par}{\par}     

% For centered blocks
\newenvironment<>{varblock}[2][.9\textwidth]{%
    \setlength{\textwidth}{#1}
    \begin{actionenv}#3%
        \def\insertblocktitle{#2}%
        \par%
        \usebeamertemplate{block begin}}
    {\par%
        \usebeamertemplate{block end}%
    \end{actionenv}}

% For table of contents centering
\newsavebox{\longestsec}% Box to save longest sectional heading

% For removing space after equations
\makeatletter
\g@addto@macro\normalsize{%
  \setlength\abovedisplayskip{0.2cm}
  % \setlength\belowdisplayskip{-0pt}
  \setlength\belowdisplayskip{0.2cm}
}

\usepackage{graphicx} % For adding images to the document
\graphicspath{{../}} % Set the parent path for images

% \usepackage{xcolor}
% \PassOptionsToPackage{dvipsnames}{xcolor} % Load other colors
% \usepackage{subfigure}
% \usepackage{caption}
\usepackage[labelfont=bf, justification=justified, skip=0pt]{caption}
\usepackage[labelformat=empty, justification=centering, skip=0pt]{subcaption}

\usepackage{booktabs} % For better looking tables
\usepackage{tabu} % For full width tables
\usepackage{makecell} % For multiline cells in tables

% Disegni LaTeX
\usepackage{tikz}
\usetikzlibrary{shapes}
%\usetikzlibrary{shapes,snakes}
%\usetikzlibrary{trees}
\newcommand{\ImageWidth}{11cm}
\usetikzlibrary{decorations.pathreplacing, positioning, arrows.meta}
\usepackage{pgfplots,siunitx}
\usepackage{adjustbox}

\pgfplotsset{compat=1.9}
\usepgfplotslibrary{units}

%%%%%%%%%%%%%%%%%%%%%%%%%%%%%%%%%%%%%%%%%%%%%%%%%%%%%%%%%%%%%%%%%%%%%%%%%%%%%%%%%%%%%%%%%%%%%%%%%%%%%

%Nuovi comandi
\newcommand{\magnitude}[1]{\left\lVert #1 \right\rVert}
\DeclareMathOperator*{\argmin}{arg\,min}  % Declaring operator argmin
\DeclareMathOperator*{\argmax}{arg\,max}  % Declaring operator argmax

% To make equal with 'def' above in equations
\newcommand\eqdef{\mathrel{\overset{\makebox[0pt]{\mbox{\normalfont\tiny def}}}{=}}}

% For the norm in equations
\newcommand\norm[1]{\left\lVert#1\right\rVert}

% Stars made from fontawesome, xcolor and stackengine
\newcommand{\mystar}{\bottominset{\textcolor{Dandelion}{\faStarO}}{\textcolor{Goldenrod}{\faStar}}{}{}}
\newcommand{\mystarO}{\textcolor{Dandelion}{\faStarO}}


%%%%%%%%%%%%%%%%%%%%%%%%%%%%%%%%%%%%%%%%%%%%%%%%%%%%%%%%%%%%%
\usetheme{Pittsburgh}
% \usetheme{CambridgeUS}
%\mode<presentation>{\usetheme{Boadilla}}
%singapore

\definecolor{amber(sae/ece)}{rgb}{1.0, 0.49, 0.0}
\definecolor{ao(english)}{rgb}{0.0, 0.5, 0.0}
\definecolor{myblue}{RGB}{0, 82, 155}
\definecolor{myred}{RGB}{140, 8, 8}

%\useoutertheme[right]{sidebar}
\usecolortheme{seahorse}
% \usecolortheme{beaver}
%\usecolortheme{dolphin}
%\definecolor{UBCblue}{rgb}{0.04706, 0.13725, 0.26667} % UBC Blue (primary)
%\usecolortheme[named=UBCblue]{structure}
%\usecolortheme[named=UBCblue]{seahorse}
%\usecolortheme{structure}
% \usecolortheme[RGB={191, 15, 15}]{structure} 
% \usecolortheme[RGB={140, 8, 8}]{structure} 
\usecolortheme[named=myred]{structure}
% \useoutertheme{modificatema2}
\setbeamertemplate{navigation symbols}{} % Remove navigation buttons
% \setbeamertemplate{footline}[frame number] % Add page number

\usefonttheme{professionalfonts}
\setbeamercovered{dynamic}

\theoremstyle{definition}
\newtheorem{definizione}{Definizione}[section]
%\theoremstyle{plain}
%\newtheorem{teorema}{Teorema}


\title{Application of GLM Advancements \\ to Non-Life Insurance Pricing}
%\subtitle{}
\author{Leonardo Stincone}

\date{18 Maggio 2021}
\institute[units]{Università degli Studi di Trieste}

\logo{\includegraphics[width=10mm, height=10mm]{figures/university-logo.png}}


% % For table of contents positioning
% % \newsavebox{\longestsec}% Box to save longest sectional heading
% \usepackage{varwidth}    
% \usepackage{etoolbox}
% \makeatletter
% %\patchcmd{\beamer@sectionintoc}{\vskip1.5em}{\vskip0.5em}{}{}
% \patchcmd{\beamer@sectionintoc}{%
%   \hbox{\vbox{%
%     \def\beamer@breakhere{\\}%
%     \beamer@tocact{\ifnum\c@section=#1\beamer@toc@cs\else\beamer@toc@os\fi}{section in toc}}}%
% }{%
%   \hbox{%
%     \def\beamer@breakhere{}%
%     \beamer@tocact{\ifnum\c@section=#1\beamer@toc@cs\else\beamer@toc@os\fi}{section in toc}}%
% }{}{}
% \makeatother    

% \newcommand\Closer[2][1cm]{%
% \makebox[\linewidth][c]{%
%   \begin{minipage}{\dimexpr\textwidth+#1\relax}
%   \raggedright#2
%   \end{minipage}%
%   }%
% }


\setbeamertemplate{section in toc}[sections numbered]
% \setbeamertemplate{section in toc}



% Insert table of content at the beginning of each section
% \AtBeginSubsection[]
% {
  % \begin{frame}[c]%{Indice}
    % \frametitle{Indice}
    % \tableofcontents[currentsection, currentsubsection]
  % \end{frame}
% }



%%%%%%%%%%%%%%%%%%%%%%%%%%%%%%%%%%%%%%%%%%%%%%%%%%%%%%%%%%%%%%%%%%%%%%%%%%%%%%%%%%%%%%%%%%%%%%%%%%%%%
%%%%%%%%%%%%%%%%%%%%%%%%%%%%%%%%%%%%%%%%%%%%%%%%%%%%%%%%%%%%%
\IfFileExists{upquote.sty}{\usepackage{upquote}}{}
\begin{document}


%% title frame
\begin{frame}
\titlepage
\end{frame}




%%%%%%%%%%%%%%%%%%%%%%%%%%%%%%%%%%%%%%%%%%%%%%%%%%%%%%%%%%%%%
%%%%%%%%%%%%%%%%%%%%%%%%%%%%%%%%%%%%%%%%%%%%%%%%%%%%%%%%%%%%%
\begin{frame}
  \frametitle{Indice}
  \begin{lrbox}{\longestsec}Last section in the presentation\end{lrbox}% Capture longest title
  \setlength{\leftskip}{\dimexpr.5\textwidth-.5\wd\longestsec\relax}% Advance left margin accordingly
  \tableofcontents
\end{frame}


%%%%%%%%%%%%%%%%%%%%%%%%%%%%%%%%%%%%%%%%%%%%%%%%%%%%%%%%%%%%%
%%%%%%%%%%%%%%%%%%%%%%%%%%%%%%%%%%%%%%%%%%%%%%%%%%%%%%%%%%%%%
\section{Descrizione del problema}

%%%%%%%%%%%%%%%%%%%%%%%%%%%%%%%%%%%%%%%%%%%%%%%%%%%%%%%%%%%%%
\begin{frame}{Descrizione del problema}
  \textcolor{myred}{\textbf{Problema}}: \textbf{prevedere il numero di sinistri} ($N_i$) che causerà un assicurato ($i$) a partire dalle informazioni della sua polizza:
  % \vspace{0.2cm}
  $$ (x_{i1}, x_{i2}, \dots, x_{ip}) \ \longmapsto \ F_{N_i}, \, E(N_i), \, Var(N_i)$$
  
  % \vspace{0.4cm}
  \vspace{0.2cm}
  
  \textcolor{myred}{\textbf{Soluzione}}: stimo un \textbf{modello} a partire dai \textbf{dati storici}.
  
  \vspace{0.4cm}
  
  \textcolor{myred}{\textbf{Perché}}: prevedere il numero di sinistri è uno degli elementi per \textbf{determinare il prezzo} di una polizza assicurativa.
\end{frame}



%%%%%%%%%%%%%%%%%%%%%%%%%%%%%%%%%%%%%%%%%%%%%%%%%%%%%%%%%%%%%
%%%%%%%%%%%%%%%%%%%%%%%%%%%%%%%%%%%%%%%%%%%%%%%%%%%%%%%%%%%%%
\section{Dataset}


%%%%%%%%%%%%%%%%%%%%%%%%%%%%%%%%%%%%%%%%%%%%%%%%%%%%%%%%%%%%%
\begin{frame}
  \frametitle{Dataset: Esposizione e Variabile Risposta}
  
  \fontsize{9pt}{11pt}\selectfont
  
  \begin{center}
    \begin{minipage}{6cm}
      \begin{varblock}[6cm]{Origine del Dataset}
        Portafoglio RCA costituito da polizze di una provincia italiana nel periodo 2014-2019
      \end{varblock}
    \end{minipage}
  \end{center}
  
  \begin{table}[!h]
    \centering
    \begin{tabular}{lcccccc}
      \toprule
      \textbf{Set} & \textbf{\makecell[c]{Esposizione\\(rischi anno)}} & \textbf{\makecell[c]{Numero\\Sinistri}} & \textbf{\makecell[c]{Frequenza\\Sinistri}}\\
      \midrule[\heavyrulewidth]
      Train & 107 998.4 & 4 823 & 0.045\\
      Test  &  26 806.3 & 1 131 & 0.042\\
      \midrule
      Tot   & 134 804.7 & 5 954 & 0.044\\
      \bottomrule
    \end{tabular}
  \end{table}
\end{frame}


%%%%%%%%%%%%%%%%%%%%%%%%%%%%%%%%%%%%%%%%%%%%%%%%%%%%%%%%%%%%%
\begin{frame}
  \frametitle{Dataset: Variabili Esplicative}
  
  \fontsize{9pt}{11pt}\selectfont
  
  \begin{table}[!h]
    \centering
    \begin{tabular}{lc}
      \toprule
      \textbf{Descrizione} & \textbf{\makecell[c]{Numero di variabili\\per categoria}}\\
      \midrule[\heavyrulewidth]
      \cellcolor{gray!6}{Informazioni sul veicolo assicurato} & \cellcolor{gray!6}{12}\\
      Informazioni generiche sull'assicurato & 14\\
      \cellcolor{gray!6}{Informazioni assicurative sull'assicurato} & \cellcolor{gray!6}{9}\\
      Opzioni della polizza assicurativa & 11\\
      \cellcolor{gray!6}{Informazioni sull'assicurato in quanto cliente} & \cellcolor{gray!6}{2}\\
      Dati telematici & 4\\
      \midrule
      \textbf{Totale} & \textbf{52}\\
      \bottomrule
    \end{tabular}
  \end{table}
\end{frame}



%%%%%%%%%%%%%%%%%%%%%%%%%%%%%%%%%%%%%%%%%%%%%%%%%%%%%%%%%%%%%
%%%%%%%%%%%%%%%%%%%%%%%%%%%%%%%%%%%%%%%%%%%%%%%%%%%%%%%%%%%%%
\section{Modelli}


%%%%%%%%%%%%%%%%%%%%%%%%%%%%%%%%%%%%%%%%%%%%%%%%%%%%%%%%%%%%%
\begin{frame}
  \frametitle{Modelli: Modelli Lineari Generalizzati (GLM)}
  
  \fontsize{9pt}{11pt}\selectfont
  
  \begin{center}
    \begin{minipage}{11cm}
      % \begin{varblock}[10cm]{Origine del Dataset}
        \begin{block}{Ipotesi}
          \begin{enumerate}
            \item $(Y_1, \dots, Y_n)$ indipendenti con distribuzione appartenente alla medesima \\ Famiglia Esponenziale Lineare
            \item Predittore lineare %\\
              $$\eta_i = \beta_0 + \beta_1 x_{i1} + \beta_2 x_{i2} + \dots + \beta_p x_{ip} = \boldsymbol{x}_i^t \boldsymbol{\beta}, \quad i\in\{1,2,\dots,n\}$$
            \item Funzione legame %\\
              $$g(\mu_i) = \eta_i = \boldsymbol{x}_i^t \boldsymbol{\beta}, \quad i\in\{1,2,\dots,n\}$$
          \end{enumerate}
        \end{block}
        \begin{block}{Stima dei parametri}
          Stima di massima verosimiglianza %\\
          $$\hat{\boldsymbol{\beta}} = \argmax_{\boldsymbol{\beta}\in\mathbb{R}^{p+1}}{L\left(\boldsymbol{\beta}, \phi; \boldsymbol{y}\right)}
          = \argmin_{\boldsymbol{\beta}\in\mathbb{R}^{p+1}}{D(\boldsymbol{\beta}, \boldsymbol{y})}$$
        \end{block}
        \begin{block}{Selezione delle variabili}
          Algoritmo stepwise basato su un criterio di informazione (AIC)
        \end{block}
      % \end{varblock}
    \end{minipage}
  \end{center}
  
  % \begin{columns}[c]
  %   \begin{column}{0.65\textwidth}
  %     \begin{block}{Ipotesi}
  %       \begin{enumerate}
  %         \item $(Y_1, \dots, Y_n)$ indipendenti con distribuzione appartenente alla medesima Famiglia Esponenziale Lineare
  %         \item Predittore lineare \\
  %           $\eta_i = \beta_0 + \beta_1 x_{i1} + \beta_2 x_{i2} + \dots + \beta_p x_{ip}$
  %         \item Funzione legame \\
  %           $g(\mu_i) = \eta_i$
  %       \end{enumerate}
  %     \end{block}
  %     \begin{block}{Stima dei parametri}
  %       Stima di massima verosimiglianza
  %       $\hat{\boldsymbol{\beta}} = \argmax_{\boldsymbol{\beta}\in\mathbb{R}^{p+1}}{L\left(\boldsymbol{\beta}, \phi; \boldsymbol{y}\right)}
  %       = \argmin_{\boldsymbol{\beta}\in\mathbb{R}^{p+1}}{D(\boldsymbol{\beta}, \boldsymbol{y})}$
  %       % \begin{align*}
  %       %   \hat{\boldsymbol{\beta}}   & = \argmax_{\boldsymbol{\beta}\in\mathbb{R}^{p+1}}{L\left(\boldsymbol{\beta}, \phi; \boldsymbol{y}\right)} \\
  %       %                              & = \argmin_{\boldsymbol{\beta}\in\mathbb{R}^{p+1}}{D(\boldsymbol{\beta}, \boldsymbol{y})}
  %       %   \end{align*}
  %     \end{block}
  %   \end{column}
  %   \begin{column}{0.35\textwidth}
  %     \begin{block}{Selezione delle variabili}
  %       \begin{itemize}
  %         \item Visualizzazione
  %         \item Test di verifica di ipotesi
  %         \item Criteri di informazione \\ (AIC, BIC)
  %         \item Divisione Train/Test
  %         \item Cross validation
  %       \end{itemize}
  %     \end{block}
  %   \end{column}
  % \end{columns}
  
\end{frame}



%%%%%%%%%%%%%%%%%%%%%%%%%%%%%%%%%%%%%%%%%%%%%%%%%%%%%%%%%%%%%
\begin{frame}
  \frametitle{Modelli: Elastic Net}
  
  \fontsize{9pt}{11pt}\selectfont
  
  \begin{block}{Elastic Net}
    $$
    \hat{\boldsymbol{\beta}} = \argmin_{\boldsymbol{\beta}\in\mathbb{R}^{p+1}}{\left\{
    D(\boldsymbol{\beta}, \boldsymbol{y}) +
    \lambda 
    \sum_{j=1}^p{\left(\alpha |\beta_j| + (1 - \alpha) |\beta_j|^2\right)}
    \right\}}
    $$
    \begin{itemize}
      \item $\lambda\ge0$ iperparametro di penalizzazione
      \item $\alpha \in [0,1]$ iperparametro che determina il peso della penalizzazione LASSO
      \begin{itemize}
        \item[\faCaretRight] $\alpha=0 \ \Longrightarrow \ $Regressione Ridge
        \item[\faCaretRight] $\alpha=1 \ \Longrightarrow \ $Regressione LASSO
      \end{itemize}
    \end{itemize}
  \end{block}
  
\end{frame}


%%%%%%%%%%%%%%%%%%%%%%%%%%%%%%%%%%%%%%%%%%%%%%%%%%%%%%%%%%%%%
\begin{frame}
  \frametitle{Modelli: Modelli Additivi Generalizzati (GAM)}
  
  \fontsize{9pt}{11pt}\selectfont
  
  \begin{block}{Modelli Additivi Generalizzati (GAM)}
    Predittore lineare
    $$
    \eta_i = \boldsymbol{x}_i^t \boldsymbol{\beta} + \sum_{l=1}^{q}{f_l(z_{i,l})}, \quad i\in\{1,2,\dots,n\}
    $$
  
    $$
    \boldsymbol{\hat{f}} = \argmin_{\boldsymbol{f}}
    {\left\{
      D(\boldsymbol{f}, \boldsymbol{y})
        + \sum_{l=1}^{q}{
          \lambda_l \int_{a_l}^{b_l}{\left( f_l''(x_l) \right)^2 dx}
        }
    \right\}
    }
    $$
    con $\lambda_1, \lambda_2, \dots, \lambda_q$ iperparametri di smoothing.
  \end{block}
\end{frame}



%%%%%%%%%%%%%%%%%%%%%%%%%%%%%%%%%%%%%%%%%%%%%%%%%%%%%%%%%%%%%
%%%%%%%%%%%%%%%%%%%%%%%%%%%%%%%%%%%%%%%%%%%%%%%%%%%%%%%%%%%%%
\section{Risultati}


%%%%%%%%%%%%%%%%%%%%%%%%%%%%%%%%%%%%%%%%%%%%%%%%%%%%%%%%%%%%%
\begin{frame}
  \frametitle{Risultati}
  
  % \setlength{\abovedisplayshortskip}{-1cm}
  % \setlength{\belowdisplayshortskip}{0cm} 
  
  \fontsize{7.5pt}{9.5pt}\selectfont
  % \fontsize{9pt}{11pt}\selectfont
  
  \begin{figure}
    \centering
    \includegraphics[width=6.5cm]{_bookdown_files/_main_files/figure-latex/models-results-plot-1.pdf}
  \end{figure}
  
  \begin{table}[!h]
    % \captionsetup{skip=0pt,belowskip=0pt}
    % \caption[Models results.]{\label{tab:models-results}Models results. In the columns we reported the deviance computed in the test set, the execution time and, for the GLM-based models, the hyper-parameters.}
    \centering
    \begin{tabular}[t]{llcccc}
      \toprule
      \textbf{Id} & \textbf{Model} & \textbf{\makecell[c]{Test\\Deviance}} & \textbf{Time} & \textbf{$\alpha$} & \textbf{$\lambda$}\\
      \midrule[\heavyrulewidth]
      \cellcolor{gray!6}{Mod1} & \cellcolor{gray!6}{GLM Tot} & \cellcolor{gray!6}{8 458.147} & \cellcolor{gray!6}{2.7s} & \cellcolor{gray!6}{0} & \cellcolor{gray!6}{0}\\
      Mod2 & Elastic Net Tot & 8 458.024 & 1h 30m & 0.06 & 2.01e-04\\
      \cellcolor{gray!6}{Mod3} & \cellcolor{gray!6}{\textbf{Ridge Tot}} & \cellcolor{gray!6}{\textbf{8 457.465}} & \cellcolor{gray!6}{1h 30m} & \cellcolor{gray!6}{0} & \cellcolor{gray!6}{4.64e-04}\\
      Mod4 & GLM AIC & 8 458.023 & 7h 27m & 0 & 0\\
      \cellcolor{gray!6}{Mod5} & \cellcolor{gray!6}{Elastic Net AIC} & \cellcolor{gray!6}{8 458.236} & \cellcolor{gray!6}{8h 54m} & \cellcolor{gray!6}{0} & \cellcolor{gray!6}{1.63e-05}\\
      Mod6 & GAM AIC & 9 728.570 & 7h 45m & 0 & 0\\
      \cellcolor{gray!6}{Mod7} & \cellcolor{gray!6}{GBM Tot} & \cellcolor{gray!6}{8 504.178} & \cellcolor{gray!6}{2h 30m} & \cellcolor{gray!6}{} & \cellcolor{gray!6}{}\\
      \bottomrule
    \end{tabular}
  \end{table}
\end{frame}


%%%%%%%%%%%%%%%%%%%%%%%%%%%%%%%%%%%%%%%%%%%%%%%%%%%%%%%%%%%%%
%%%%%%%%%%%%%%%%%%%%%%%%%%%%%%%%%%%%%%%%%%%%%%%%%%%%%%%%%%%%%
\begin{frame}
  \begin{center}
    \textcolor{myred}{\Huge \textbf{Grazie per l'attenzione}}
  \end{center}
\end{frame}



\end{document}
